%!TEX root = main.tex

In the sequel, we introduce an operational semantics of LoFIRRTL programs, on which the formalization of LoFIRRTL in Coq is based. 

%At first, let us make it more precise the syntactical constraints of LoFIRRTL stated in Section~\ref{sec-firrtl}.

Let us fix a LoFIRRTL circuit $A$. Then from the syntax of LoFIRRTL, the top-level module of the circuit is a module $A$. 

We first show how to define the operational semantics for the situation that the module $A$ contains no occurrences of ``inst of'' statements. Later on, we will consider the more general case.
%
\paragraph*{No ``inst of'' statements in module $A$} Let us assume that $A$ contains no occurrences of ``inst of'' statements. In this case, we know that $A$ is the only module in circuit $A$.

Before presenting the operational semantics, we introduce some notations. 

We use $\iports(A)$ (resp. $\oports(A)$) to denote the set of input (resp. output) ports of the module $A$,  $\ports(A) = \iports(A) \cup \oports(A)$, $\wires(A)$ to denote the set of wires in the module instance $A$, $\nodes(A)$ to denote the set of nodes in the module $A$, and $\registers(A)$ to denote the set of registers in the module $A$.
%
Moreover, we use $\cpnts(A)$ to denote the set of components in $A$, that is, $\ports(A) \cup \wires(A) \cup \nodes(A) \cup \registers(A)$.

A \emph{generalized connect statement} for $z  \in \cpnts(B)$ is a statement of the form $z <= e$ or $\mbox{node } z = e$, where $z$ and $e$ are called the \emph{head} and \emph{body} of the generalized connect statement respectively. Intuitively, a generalized connect statement is either a connect statement or a node-declaration statement. 

We use $\ccpnts(A)$ to denote the set of connected components of $A$, that is, $\ccpnts(A) = \oports(A) \cup \wires(A) \cup \nodes(A)$. (Intuitively, connected components of $A$ must be connected to, that is, occur as the heads of generalized connect statements of $A$.) Note that the input ports of $A$ are not connected to, moreover, the registers may not be connected to. (If a register is not connected to, then it keeps the current value. ) 

%For $(id, B) \in \modinst(A)$, let $\cpnts(id, B)$ denote $\{id.z \mid z \in \cpnts(B)\}$.

%Let $B$ be a module in circuit $A$. A \emph{generalized connect statement} for $z  \in \cpnts(B)$ is a statement of the form $z <= e$ or $\mbox{node } z = e$, where $z$ and $e$ are called the \emph{head} and \emph{body} of the generalized connect statement respectively. (Note that in the syntax of LoFIRRTL, only statements of the form $z <= e$ are called connect statements.)

%
The \emph{dependency graph} of $A$ is $G_A = (\cpnts(A), E_A)$, where $E_A$ comprises the edges $(x, y)$ such that $y$ depends on $x$, more precisely, $x$ occurs in the body of the generalized connect statement for $y$. 
%
%(Recall that LoFIRRTL does not contain partial connect statements, moreover, for each $y \in \cpnts(B)$, there is at most one connect statement where $y$ is the head.)

Since combinational loops\footnote{A combinatorial loop is when the output of some combinatorial logic, that is a section of logic with no registers between gates, is fed back into the input of that combinatorial logic with no intervening register.} are usually taken as ``bad'' in digital circuit design, we require that \emph{$G_A$ is acyclic}. Another necessary condition for the well-definedness of semantics of LoFIRRTL is stated as follows: \emph{Each $z \in \ccpnts(A)$ should be connected to exactly once, that is, there is exactly one generalized connect statement whose head is $z$}. Futhermore, in this work, we ignore clocks, that is, all signals are related to the same clock, the global clock.

To define the operational semantics of $A$, we require that module $A$ is in \emph{pseudo-SSA form}, that is,  the statement sequence of module $A$ satisfies the following two additional  constraints,
\begin{itemize}
\item all components are declared before they are used in the bodies and heads of connect statements (except that nodes are declared and used simultaneously),
%
\item the sequence of heads of the generalized connect statements in module $A$ forms a topological sorting of $G_A$ (where the vertices without incoming edges are ignored), specifically, for every edge $(x, y)$ in $G_A$, $x$ is before $y$ in the sequence (if $x$ indeed occurs in the sequence), in other words, the generalized connect statement for $x$ (if there is any) is before the connect statement for $y$ in $A$. 
\end{itemize}

%A necessary condition for the well-definedness of the semantics of LoFIRRTL programs

%From the fact that $A$ is a module in a LoFIRRTL program, we know that $G_A$ is acyclic. 
%Equipped with $G_A$, we state a necessary condition for the well-definedness of the semantics of LoFIRRTL programs: $G_A$ satisfies the following constraints,
%\begin{itemize}
%\item for each $z \in \ccpnts(A)$, there is exactly one generalized connect statement whose head is $z$, 
%
%\item for each $z \in \registers(A)$, there is at most one generalized connect statement whose head is $z$, 
%
%\item the dependency graph $G_A$ is acyclic.
%\end{itemize}


For instance, the following LoFIRRTL program
%[caption=My FIRRTL Example]
\begin{lstlisting}
module A
  input io_in: UInt<32>
  output io_out: UInt<32>
  node a = add(io_in, UInt<32>("h1"))
  io_out <= tail(a, 1)
\end{lstlisting}
%
is in pseudo-SSA form, while the following program
\begin{lstlisting}
module A
  input io_in: UInt<32>
  output io_out: UInt<32>
  io_out <= tail(a, 1)
  node a = add(io_in, UInt<32>("h1"))
\end{lstlisting}
is not, since $({\tt a}, {\tt io\_out})$ is an edge in $G_A$, while the generalized connect statement for {\tt a} is after the generalized connect statement for {\tt io\_out}.
 
In the sequel, we require that $A$ is in the pseudo-SSA form.

The operational semantics of $A$ relies on a concept of configurations defined below. 

Let $\gtypes$ denote the set of ground types, that is, $(\{UInt, SInt\} \times \Nat) \cup \{Clock, Reset, AsyncReset\}$. Let $\bv$ denote the set of bit vectors, that is, $\{0,1\}^+$.
For $t = (\tau, n) \in \{UInt, SInt\} \times \Nat$, we define the \emph{width} of $t$, denoted by $wd(t)$, as $n$.

A \emph{configuration} of $A$ is a triple $(te, s, rs)$, where $te$ is a type environment, $s$ is a state, and $rs$ is a register state. Intuitively, $te$ denotes the types of the components, $s$ denotes the values of the components in the beginning of the current clock cycle, and $rs$ denotes the new values of the registers in the beginning of the next clock cycle. Formally, 
\begin{itemize}
\item a \emph{type environment} $te$ of $A$ is a function from $\cpnts(A)$ to $\gtypes$,
%
\item a \emph{state} $s$ of $A$ is a function from $\cpnts(A)$ to $\bv \cup \{\bot\}$,
%
\item a \emph{register state} $rs$ of $A$ is a function from $\registers(A)$ to $\bv \cup \{\bot\}$.
\end{itemize}

%A \emph{type environment} of $M$ is a function from $\varset(M)$ to $\gtypes$.
%A \emph{state} of $M$ is a function from $\varset(M)$ to $\bv \cup \{\bot\}$. A \emph{register state} $rs$ of $M$ is a function from $\varset(M)$ to $\bv \cup \{\bot\}$ such that $rs(v) = \bot$ for every $v \in \gtypes \setminus \registers(M)$. Note that $\bot$ is a special symbol denoting the fact that the value of a variable is not determined yet. 

The operational semantics of each statement $st$ in $A$ is defined as a relation $(te, s, rs) \xrightarrow{st} (te', s', rs')$, which denotes that the execution of $st$ from a configuration $(te, s, rs)$ results into a configuration $(te', s', rs')$. Moreover, the operational semantics of a sequence of statements is defined as the composition of the relations of statements therein and the operational semantics of module $A$ is defined as the operational semantics of its statement sequence.


%Then a \emph{configuration} of $M$ is defined as a triple $(te, s, rs)$, where $te$ is a type environment, $s$ is a state, and $rs$ is a register state. Intuitively, $te$ denotes the types of the variables, $s$ denotes the values of the variables in the beginning of the current clock cycle, and $rs$ denotes the new values of the registers in the beginning of the next clock cycle. 

%Then the operational semantics of each statement $st$ in $M$ is defined as a relation $(te, s, rs) \xrightarrow{st} (te', s', rs')$, which denotes that the execution of $st$ from a configuration $(te, s, rs)$ results into a configuration $(te', s', rs')$. Moreover, the operational semantics of a sequence of statements is defined as the composition of the relations of statements therein. 


The rules for the operational semantics of statements are presented in Table~\ref{tab:lofirrtl-sem}, where $te[t/v]$ denotes the function that is obtained from $te$ by assigning $t$ to $v$, similarly for $s[0^{wd(t)}/v]$, and so on. Moreover, the rules rely on $te(e)$ and $s(e)$, which are extensions of the functions $te$ and $s$ to expressions, whose definitions can be found in Appendix~\ref{app-lofirrtl}. 

Let us illustrate the semantics by using the ``reg with reset: $\neq \bot$, true'' rule: At first, $rs(r) \neq \bot$ implies that the value of the register $r$ in the beginning of the current clock cycle is defined. Then the state $s$ is updated by assigning $rs(r)$ to $r$, resulting into the new state $s'$. Moreover, because $s(e_1)=true$, the register state $rs$ is updated by assigning $s(e_2)$ to $r$, resulting into the new register state $rs'$. In other words, the value of $r$ is reset, that is, the value of $r$ will be $s(e_2)$ in the beginning of the next clock cycle.
%

\begin{table}[t]
  \footnotesize
  \begin{gather*}
    \infer[skip]
    {\seqq{(te, s, rs) \xrightarrow{\mbox{skip}} (te, s, rs)}}
    {\seqq{-}}
    \quad
    \inferii[wire]
%    {\seqq{}}
    {\seqq{(te, s, rs) \xrightarrow{\mbox{wire } v: t} (te', s', rs)}}
   {\seqq{te' = te [t/v]}} {\seqq{s' = s [\bot/v]}}
    \\[1ex]
    \inferiii[reg: $\bot$]
    {\seqq{(te, s, rs) \xrightarrow{\mbox{reg } r: t, c} (te', s', rs)}}
   {\seqq{rs(r) = \bot}} {\seqq{te' = te [t/r]}} {\seqq{s' = s[\bot/r]}}
    \quad
    \inferii[reg: $\neq\bot$]
    {\seqq{(te, s, rs) \xrightarrow{\mbox{reg } r: t, c} (te, s', rs)}}
   {\seqq{rs(r) \neq \bot}} {\seqq{s' = s[rs(r)/r]}}
    \\[1ex]
    \inferiv[reg with reset: $\bot$, false]
    {\seqq{(te, s, rs) \xrightarrow{\mbox{reg } r: t, c \mbox{ with: reset => }(e_1, e_2)} (te', s', rs)}}
   {\seqq{rs(r) = \bot}} {\seqq{ s(e_1)= false}}  {\seqq{te' = te [t/r]}} {\seqq{s' = s[\bot/r]}}
    \\[1ex]
    \inferv[reg with reset: $\bot$, true]
    {\seqq{(te, s, rs) \xrightarrow{\mbox{reg } r: t, c \mbox{ with: reset => }(e_1, e_2)} (te', s', rs')}}
   {\seqq{rs(r) = \bot}} {\seqq{ s(e_1)= true}}  {\seqq{te' = te [t/r]}} {\seqq{s' = s[\bot/r]}} {\seqq{rs' = s[s(e_2)/r]}}
    \\[1ex]
    \inferiii[reg with reset: $\neq\bot$, false]
    {\seqq{(te, s, rs) \xrightarrow{\mbox{reg } r: t, c \mbox{ with: reset => }(e_1, e_2)} (te, s', rs)}}
   {\seqq{rs(r) \neq \bot}} {\seqq{ s(e_1)= false}} {\seqq{s' = s[rs(r)/r]}}
    \\[1ex]
    \inferiv[reg with reset: $\neq\bot$, true]
    {\seqq{(te, s, rs) \xrightarrow{\mbox{reg } r: t, c \mbox{ with: reset => }(e_1, e_2)} (te, s', rs')}}
   {\seqq{rs(r) \neq \bot}} {\seqq{ s(e_1)= true}} {\seqq{s' = s[rs(r)/r]}} {\seqq{rs' = rs[s(e_2)/r]}}
    \\[1ex]
    \inferii[node]
    {\seqq{(te, s, rs) \xrightarrow{\mbox{node } v = e} (te', s', rs)}}
   {\seqq{te' = te[te(e)/v]}} {\seqq{s' = s[s(e)/v]}}
    \quad
    \inferii[connect: non-register]
    {\seqq{(te, s, rs) \xrightarrow{v <= e} (te, s', rs)}}
   {\seqq{v \not \in \registers(A)}} {\seqq{s' = s[s(e)/v]}}
    \\[1ex]    
    \inferii[connect: register]
    {\seqq{(te, s, rs) \xrightarrow{v <= e} (te, s, rs')}}
   {\seqq{v \in \registers(A)}} {\seqq{rs' = rs[s(e)/v]}}
  \quad
    \infer[invalid]
    {\seqq{(te, s, rs) \xrightarrow{v \mbox{ is invalid}} (te, s', rs)}}
%   {\seqq{rs(v) \neq \bot}} {\seqq{rs' = rs[s(e)/v]}}
   {\seqq{s' = s[0^{wd(te(v))}/v]}}
  \end{gather*}
%  \vspace{-4mm}
  \caption{Operational semantics of LoFIRRTL statements.}
  \label{tab:lofirrtl-sem}
  \vspace{-6mm}
\end{table}

%$te(e)$ is defined as follows. 

An \emph{input trace} of $A$ of length $m$ is a function $it$ from $\iports(A)$ to $\bv$ such that $it(x)$ is a sequence of bit vectors and the length of $it(x)$ is $m$, for every $x \in \iports(A)$.

Let $st$ be the sequence of statements in $A$ and $it$ be an input trace of $A$ of length $m$. Then the semantics of $A$ with respect to $it$ is defined by the relations $(te, s, rs) \xrightarrow{A, it, j, m} (te', s', rs')$ for $0 \le j \le m$. Intuitively, the relation $(te, s, rs) \xrightarrow{A, it, j, m} (te', s', rs')$ denotes that $(te, s, rs)$ is the configuration resulted from processing the first $m-j$ clock cycles of the input trace and $(te', s', rs')$ is obtained after processing the remaining $j$ clock cycles of the input trace. In the sequel, we define the relation $(te, s, rs) \xrightarrow{A, it, j, m} (te', s', rs')$ inductively as follows.
\begin{itemize}
\item At first, we read the values of the input ports in the $(m-j+1)$-th clock cycle of $it$, that is, replace the value of each input port $x$ in $s$ with $it(x)[m-j+1]$, resulting into a state $s_1$.
%
\item Then, we execute the statement sequence $st$ from the configuration $(te, s_1, rs)$, resulting into a configuration $(te'', s'', rs'')$. 
%
\item Finally, we obtain $(te', s', rs')$ by processing the last $j-1$ clock cycles of the input trace $it$.
\end{itemize}
More formally, $(te, s, rs) \xrightarrow{A, it, j, m} (te', s', rs')$ is defined as follows. 
\begin{itemize}
\item If $j = 0$, then $(te, s, rs) \xrightarrow{A, it, j, m} (te, s, rs)$. 
%
\item Otherwise, let $s_1$ be obtained from $s$ by assigning $it(x)[m-j+1]$ to $x$ for each $x \in \iports(A)$. If $(te, s_1, rs) \xrightarrow{st} (te'', s'', rs'')$ and $(te'', s'', rs'') \xrightarrow{A, it, j-1, m} (te', s', rs')$, then we have $(te, s, rs) \xrightarrow{A, it, j, m} (te', s', rs')$.
\end{itemize}
%$s(e)$ is defined as follows.
%\begin{itemize}
%\item 
%\end{itemize} 

\paragraph*{There exist ``inst of'' statements in module $A$}
In this case, circuit $A$ contains at least two instances of modules. Moreover, there are no mutual dependencies between modules, that is, the module dependency graph, which comprises the edges $(B, C)$ such that a statement \begin{lstinline}! inst y of C!\end{lstinline} occurs in module $B$ for some $y$, is acyclic. 
Our main idea is to flatten the hierarchies of modules in circuit $A$ into a circuit containing a single module $M_A$, then transform $M_A$ into pseudo-SSA form $M'_A$, and finally define the operational semantics of $A$ as that of $M'_A$. (Recall that the semantics of a module that is in pseudo-SSA form and contains no occurrences of ``inst of'' statements is defined by the rules in Table~\ref{tab:lofirrtl-sem}.)

In the sequel, we show how to flatten the hierarchies of modules in circuit $A$ into a module $M_A$.

We define the set of module instances of circuit $A$, denoted by $\modinst(A)$, inductively as follows. An element of $\modinst(A)$ is a pair $(id, B)$, which denotes an instance of module $B$, with the identifier $id$. 
% \emph{the identifiers} of these module instances inductively as follows. 
\begin{itemize}
\item $(A, A) \in \modinst(A)$, that is, an instance of the top-level module $A$ is in $MI(A)$. Note that we use $A$ as the identifier of the instance of module $A$, since there is a unique instance of the top-level module.  
%
\item Suppose $(x, B) \in \modinst(A)$, and the module $B$ contains a statement \begin{lstinline}! inst y of C!\end{lstinline}. Then $(x.y, C) \in \modinst(A)$.
\end{itemize}

For each $(id, B) \in MI(A)$, we inductively construct a module $M_{(id, B)}$ as follows. 
\begin{itemize}
\item If $B$ contains no occurrences of ``inst of'' statements, then $M_{(id, B)}$ is obtained from module $B$ by replacing all identifiers $z$ with $id.z$.
%
\item Otherwise, suppose that $M_{(id.y, C)}$ has been constructed for each statement \begin{lstinline}! inst y of C!\end{lstinline} in module $B$. Then $M_{(id, B)}$ is obtained from module $B$ by first replacing all identifiers $z$ with $id.z$,  then replacing each statement \begin{lstinline}! inst id.y of C!\end{lstinline} with the statement sequence in $M_{(id.y, C)}$.
\end{itemize}
Finally, let $M_A:=M_{A, A}$.

Furthermore, $M_A$ can be turned into $M'_A$ in pseudo-SSA form by constructing the dependency graph $G_{M_A}$, building a topological sorting of $G_{M_A}$, and reordering the generalized connect statements of $M_A$ according to the topological sorting.


%%%%%%%%%%%%%%%%%%%%%%%%%%%%%%%%%%%%%%%%%%%%
%%% the original non-flattened version removed
%%%%%%%%%%%%%%%%%%%%%%%%%%%%%%%%%%%%%%%%%%%%

\hide{
Equipped with $MI(A)$ and extended dependency graphs, we inductively construct for each $(id, B) \in MI(A)$ a flattened statement sequence, denoted by $fss_{(id, B)}$, as follows. 
\begin{itemize}
\item If $B$ contains no occurrences of ``inst of'' statements, then $fss_{(id, B)}$ is constructed by the following procedure.
\begin{itemize}
\item At first, reorder the sequence of statements in $B$ into $ss'$, such that the sequence of heads of the generalized connect statements in $ss'$ forms a topological sorting of $EG_B$. 
%
\item Then, replace each identifier $z$ in $ss'$ by $id.z$.
%
\item Let $fss_{(id, B)}:=ss'$.
\end{itemize}


Let $B$ be a module. If $B$ contains a statement \begin{lstinline}! inst y of C!\end{lstinline}, then the input and output ports of $y$ are visible to $B$. 
%
We use $\siports(B)$ (resp. $\soports(B)$) to denote the set of $y.p$ such that \begin{lstinline}! inst y of C!\end{lstinline} is a  statement in module $B$ and $p \in \iports(C)$ (resp. $p \in \oports(C)$). Let $\sports(B) = \siports(B) \cup \soports(B)$. Moreover, we define $\cpnts(B) = \ports(B) \cup \wires(B) \cup \nodes(B) \cup \registers(B) \cup \sports(B)$ and $\ccpnts(B) = \oports(B) \cup \wires(B) \cup \nodes(B) \cup \siports(B)$.

Let $B$ be a module. We define the \emph{extended dependency graph} of $B$, denoted by $EG_B = (V_B, E_B)$, as follows. $V_B = \cpnts(B)$ and $E_B$ is defined in the sequel. 
\begin{itemize}
\item If $B$ contains no occurrences of ``inst of'' statements, then $E_B$ comprises the edges $(x, y)$ such that $y$ depends on $x$, that is, there is a connect statement in $B$ where $y$ is the head and $x$ occurs in the body. 
%
\item Otherwise, suppose that $M_{(id.y, C)}$ has been constructed for each statement \begin{lstinline}! inst y of C!\end{lstinline} in module $B$. Then 
%suppose that the graph $EG_C$ has been constructed for each statement \begin{lstinline}! inst y of C!\end{lstinline} in module $B$, 
%then 
$E_B$ comprises the following two types of edges, 
\begin{itemize}
\item the edges $(x, y)$ such that there is a connect statement in $B$ where $y$ is the head and $x$ occurs in the body,  
%
\item the edges $(y.p_1, y.p_2)$ such that \begin{lstinline}! inst y of C!\end{lstinline} is a  statement in module $B$, $p_1 \in \iports(C)$ and $p_2 \in \oports(C)$.
%, moreover, $p_2$ is reachable from $p_1$ in $EG_C$.
\end{itemize}
Note here we add to $E_B$ the edges from every input port of $C$ to every output port of $C$.
\end{itemize}

Equipped with $MI(A)$ and extended dependency graphs, we inductively construct for each $(id, B) \in MI(A)$ a flattened statement sequence, denoted by $fss_{(id, B)}$, as follows. 
\begin{itemize}
\item If $B$ contains no occurrences of ``inst of'' statements, then $fss_{(id, B)}$ is constructed by the following procedure.
\begin{itemize}
\item At first, reorder the sequence of statements in $B$ into $ss'$, such that the sequence of heads of the generalized connect statements in $ss'$ forms a topological sorting of $EG_B$. 
%
\item Then, replace each identifier $z$ in $ss'$ by $id.z$.
%
\item Let $fss_{(id, B)}:=ss'$.
\end{itemize}
%
\item Otherwise, suppose that $fss_{(id.y, C)}$ has been constructed for each statement \begin{lstinline}! inst y of C!\end{lstinline} in module $B$. Then $fss_{(id, B)}$ is constructed by the following procedure. 
\begin{itemize}
\item We first construct a topological sorting of $EG_B$, say $z_1, \cdots, z_m$. For technical convenience, we assume that the input ports of $B$ are ignored in $z_1,\cdots, z_m$. 
%
\item Let $ss'$ be an enumeration of the declaration statements for ports, wires, and registers in $B$.  
For each $y$ such that $B$ contains a statement of the form \begin{lstinline}! inst y of C!\end{lstinline}, define $idx_y$ as the minimum index $j \in [m]$ such that $z_j = y.p$ for some output port $p$ of $C$. 
Then let $i:=1$ and iterate the following procedure until $i=m$.
\begin{itemize}
\item If $z_i \not \in \soports(B)$, let $st_i$ be the connect statement of $B$ where $z_i$ is the head and $st'_i$ be obtained from $st_i$ by replacing each identifier $z'$ with $id.z'$, then $ss' := ss' \concat st'_i$. 
\item Otherwise, $z_i \in \soports(B)$. Suppose $z_i = y.p$ such that \begin{lstinline}! inst y of C!\end{lstinline} is a statement in $B$. If $i = idx_y$, then let $ss':=ss' \concat fss_{(id.y, C)}$.
\end{itemize}
%
\item Let $fss_{(id, B)}:= ss'$.
\end{itemize}
\end{itemize}
}
%%%%%%%%%%%%%%%%%%%%%%%%%%%%%%%%%%%%%%%%%%%%
%%% the original non-flattened version removed
%%%%%%%%%%%%%%%%%%%%%%%%%%%%%%%%%%%%%%%%%%%%
